\section{Generalized OT Problems}
\label{sec-generalized-ot-problems}
\subsection{OT Barycenters}
\label{sec-barycenters}
\paragraph{Fr\'echet means.}
The natural formulation is a Fr\'echet-mean problem on the space of probability measures: the unknown is the barycenter measure itself, and its support is not prescribed. It uses the continuous Kantorovich value $\MK_c$ defined in~\eqref{eq-mk-generic}.

\begin{defn}[Optimal-transport barycenter]\label{def-ot-barycenter}
Given input measures $(\be_s)_{s=1}^S$ on a space $\Xx$ and weights $\la\in\simplex_S$, discard any zero-weight inputs. An \emph{optimal-transport barycenter} of this weighted family is any solution of
\eql{\label{eq-barycenter-generic}
\umin{\al \in \Mm_+^1(\Xx)} \sum_{s=1}^S \la_s \MK_{\c}(\al,\be_s).
}
\end{defn}
\paragraph{Fixed-support restriction.}
\(\be_s=\sum_{j=1}^{n_s} b_{s,j}\de_{x_{s,j}}, \qquad \b_s=(b_{s,j})_{j=1}^{n_s}\in\simplex_{n_s}.\)
\((\C_s)_{i,j}=c(y_i,x_{s,j}) \in \RR^{n\times n_s}.\)
Thus its value is $\MKD_{\C_s}(\a,\b_s)$, in the notation of the discrete Kantorovich problem~\eqref{eq-kanto-discr}.

\begin{defn}[Fixed-support discrete OT barycenter]\label{def-fixed-support-discrete-barycenter}
With the candidate sites and cost matrices above, a \emph{fixed-support discrete OT barycenter} is a measure $\al^\star=\sum_i a_i^\star\de_{y_i}$ whose weight vector $\a^\star=(a_i^\star)_{i=1}^n$ solves
\eql{\label{eq-wass-discr}
\umin{\a \in \simplex_n} \sum_{s=1}^S \la_s \MKD_{\C_s}(\a,\b_s),
}
\end{defn}
\[
\umin{\substack{\a\in\simplex_n\\Y=(y_i)_{i=1}^n\in(\RR^d)^n}}
\sum_{s=1}^S\la_s\MKD_{\C_s(Y)}(\a,\b_s),
\qquad
[\C_s(Y)]_{i,j}=\norm{y_i-x_{s,j}}^2.
\]
For fixed $Y$, the objective is the convex problem~\eqref{eq-wass-discr} in $\a$; its dependence on $Y$ is generally nonconvex, so the joint problem is nonconvex. Thus choosing the support is itself a significant computational part of the barycenter problem, rather than a harmless preprocessing step.

\begin{prop}[Two-measure barycenters are Wasserstein geodesics]\label{prop-two-measure-barycenter-geodesic}
Let $\be_0,\be_1\in\Pp_2(\RR^d)$ and $t\in[0,1]$. For $S=2$, $c(x,y)=\norm{x-y}^2$ and weights $(1-t,t)$, the minimizers of the barycenter problem~\eqref{eq-barycenter-generic} are exactly the time-$t$ points of constant-speed $\Wass_2$ geodesics from $\be_0$ to $\be_1$. The minimum value is
\(t(1-t)\Wass_2^2(\be_0,\be_1).\)
In particular, if $\pi$ is any optimal coupling between $\be_0$ and $\be_1$, then
\(\al_t \eqdef \bigl((x,y)\mapsto(1-t)x+ty\bigr)_\sharp\pi\)
is a barycenter. If $\be_0$ has a density, the barycenter is unique and
\(\al_t=((1-t)\Id+tT)_\sharp\be_0,\)
where $T$ is the Brenier map from $\be_0$ to $\be_1$.
\end{prop}
\begin{proof}
The endpoint cases $t\in\{0,1\}$ are immediate, so assume $0<t<1$. For any candidate $\al\in\Pp_2(\RR^d)$, set
\(A=\Wass_2(\be_0,\al), \qquad B=\Wass_2(\al,\be_1), \qquad D=\Wass_2(\be_0,\be_1).\)
The identity
\((1-t)A^2+tB^2 = t(1-t)(A+B)^2+\bigl((1-t)A-tB\bigr)^2\)
and the triangle inequality $A+B\geq D$ give
\((1-t)\Wass_2^2(\al,\be_0) +t\Wass_2^2(\al,\be_1) \geq t(1-t)D^2.\)
For the interpolation $\al_t$ induced by an optimal coupling, Proposition~\ref{prop-plan-interpolation-w2-geodesic} gives $A=tD$ and $B=(1-t)D$, so the lower bound is attained.

Conversely, every minimizer must attain equality in both inequalities above. Hence $A+B=D$ and $(1-t)A=tB$, which imply $A=tD$ and $B=(1-t)D$. Parameterize a constant-speed geodesic from $\be_0$ to $\al$ on $[0,t]$ and one from $\al$ to $\be_1$ on $[t,1]$; both pieces then have speed $D$. Their concatenation is globally geodesic: any shorter path between points on opposite pieces, combined with the remaining subsegments, would contradict $D=\Wass_2(\be_0,\be_1)$. Thus the concatenated geodesic passes through $\al$ at time $t$. Finally, if $\be_0$ has a density, Brenier's theorem (Theorem~\ref{thm-brenier}) gives the unique optimal coupling $(\Id,T)_\sharp\be_0$, and the displayed McCann interpolation follows from Definition~\ref{def-monge-mccann-interpolation}.
\end{proof}
\begin{prop}[Mean and support of quadratic Wasserstein barycenters]\label{prop-w2-barycenter-mean-support}
Let $\be_1,\ldots,\be_S\in\Pp_2(\RR^d)$, let $\la\in\simplex_S$, and let $\al^\star$ be a barycenter for the squared Euclidean cost. Define
\[
K
\eqdef
\overline{\operatorname{conv}}
\left(
\bigcup_{\{s:\,\la_s>0\}}\supp(\be_s)
\right).
\]
Then
\(\supp(\al^\star)\subset K, \qquad \int_{\RR^d}x\,\d\al^\star(x) = \sum_{s=1}^S\la_s\int_{\RR^d}x\,\d\be_s(x).\)
\end{prop}
\begin{proof}
Write $m_\al=\int x\,\d\al(x)$ and let $\al^0$ be its centered translate, as in Proposition~\ref{prop-wasserstein-mean-decomposition}. That proposition gives, for every candidate $\al$,
\(\sum_{s=1}^S\la_s\Wass_2^2(\al,\be_s) = \sum_{s=1}^S\la_s\norm{m_\al-m_{\be_s}}^2 +\sum_{s=1}^S\la_s\Wass_2^2(\al^0,\be_s^0).\)
For a fixed centered law $\al^0$, the second sum is independent of $m_\al$, whereas the first is uniquely minimized at
$m_\al=\sum_s\la_s m_{\be_s}$. Translating $\al^\star$ without changing its centered law therefore proves the mean identity.

For the support statement, let $P_K$ be the Euclidean metric projection onto the closed convex set $K$. For every $y\in K$, the projection inequality gives \(\norm{x-y}^2 \geq \norm{P_K(x)-y}^2+\norm{x-P_K(x)}^2.\)
For each $s$ with $\la_s>0$, let $\pi_s$ be an optimal coupling between $\al^\star$ and $\be_s$. Since $\be_s(K)=1$, pushing $\pi_s$ forward by $(P_K,\Id)$ and integrating the preceding inequality yields
\(\Wass_2^2\bigl((P_K)_\sharp\al^\star,\be_s\bigr) \leq \Wass_2^2(\al^\star,\be_s) - \int_{\RR^d}\operatorname{dist}(x,K)^2\,\d\al^\star(x).\)
After multiplication by $\la_s$ and summation, barycenter optimality forces the last integral to vanish. Hence $\al^\star(K)=1$, and the closedness of $K$ gives $\supp(\al^\star)\subset K$.
\end{proof}
If all input supports are compact, their finite union is compact and so is its convex hull; in that case the closure in the definition of $K$ is unnecessary.

\paragraph{One-dimensional case.}
\begin{prop}[Quantile barycenters on the line]\label{prop-quantile-barycenters}
For $\Xx=\RR$ and $c(x,y)=|x-y|^2$, the quantile function of a Wasserstein barycenter is the weighted average of the input quantile functions:
\(\cumul{\al^\star}^{-1}(r) = \sum_{s=1}^S\la_s\cumul{\be_s}^{-1}(r), \qquad r\in(0,1).\)
\end{prop}
\begin{proof}
The one-dimensional formula~\eqref{eq-wass-cumul} gives \(\sum_s\la_s\Wass_2^2(\al,\be_s) = \int_0^1 \sum_s\la_s \abs{\cumul{\al}^{-1}(r)-\cumul{\be_s}^{-1}(r)}^2 \d r.\)
The minimization decouples pointwise in $r$. For each fixed $r$, the minimizer of $z\mapsto\sum_s\la_s|z-\cumul{\be_s}^{-1}(r)|^2$ is the weighted average $\sum_s\la_s\cumul{\be_s}^{-1}(r)$. This function is nondecreasing because it is a positive weighted sum of nondecreasing quantile functions, hence it is a valid quantile function.
\end{proof}
\paragraph{Gaussian case.}
\begin{prop}[Nondegenerate Gaussian inputs remain Gaussian]\label{prop-gaussian-barycenter}
Let the positive-weight inputs be $\be_s=\Gaussian(\mean_s,\cov_s)$, with $\cov_s\succeq0$, and assume that at least one $\cov_s$ is positive definite. The quadratic Wasserstein barycenter is unique and has the form $\al^\star=\Gaussian(\mean,\cov)$, where
\(\mean=\sum_s\lambda_s\mean_s.\)
Its positive-definite covariance $\cov$ is the unique minimizer of the Bures objective
\eql{\label{eq-gaussian-barycenter-energy}
\mathcal E(X)
\eqdef
\sum_s \la_s \Bb(X,\cov_s)^2.
}
For $X\in\mathbb S_{++}^d$, define
\begin{align}
R_s(X)
&\eqdef
\pa{X^{1/2}\cov_sX^{1/2}}^{1/2},
&
T_s(X)
&\eqdef
X^{-1/2}R_s(X)X^{-1/2},
\label{eq-gaussian-barycenter-maps}
\\
\overline T(X)
&\eqdef
\sum_s\la_sT_s(X),
&
\Psi_{\lambda}(X)
&\eqdef
\sum_s\la_sR_s(X).
\label{eq-gaussian-barycenter-average-map}
\end{align}
Here $T_s(X)$ is the positive-semidefinite linear optimal map from $\Gaussian(0,X)$ to $\Gaussian(0,\cov_s)$. The covariance $\cov$ is equivalently characterized by either of the stationarity equations
\eql{\label{eq-gaussian-barycenter-fixed-point}
\overline T(\cov)=\Id
\qquad\Longleftrightarrow\qquad
\Psi_{\lambda}(\cov)=\cov.
}
In dimension one, if $\sigma_s=\sqrt{\cov_s}$ denotes the input standard deviation, then the barycenter standard deviation is $\sigma=\sum_s\lambda_s\sigma_s$.
\end{prop}
\begin{proof}
Let $\GaussProj(\al)$ denote the Gaussian measure with the same mean and covariance as $\al$. For every competitor $\al\in\Pp_2(\RR^d)$ and every Gaussian input $\be_s$, the Gelbrich contraction of Theorem~\ref{thm-gelbrich-projection} gives
\(\Wass_2^2\bigl(\GaussProj(\al),\be_s\bigr) = \Wass_2^2\bigl(\GaussProj(\al),\GaussProj(\be_s)\bigr) \leq \Wass_2^2(\al,\be_s).\)
Summing with weights $\lambda_s$ shows that moment-matched Gaussian projection cannot increase the barycenter objective. Since a barycenter exists, projecting any minimizer therefore produces a Gaussian barycenter. The input with positive-definite covariance is absolutely continuous; the uniqueness criterion stated after Definition~\ref{def-ot-barycenter} then implies that the barycenter itself is this Gaussian measure.

For a Gaussian candidate, formula~\eqref{eq-dist-gauss} separates the objective as
\[
\sum_s\lambda_s\Wass_2^2\bigl(\Gaussian(\mean,\cov),\Gaussian(\mean_s,\cov_s)\bigr)
=
\sum_s\lambda_s\norm{\mean-\mean_s}^2
+
\sum_s\lambda_s\Bb(\cov,\cov_s)^2.
\]
The first term is uniquely minimized at $\mean=\sum_s\lambda_s\mean_s$. The second is the Bures barycenter problem. Uniqueness of the Wasserstein barycenter makes its covariance minimizer unique, and the presence of a positive-definite input makes this minimizer positive definite.

First suppose that $\cov_s\succ0$. Cyclicity of the trace gives \(\tr R_s(X) = \tr\pa{\cov_s^{1/2}X\cov_s^{1/2}}^{1/2}.\)
For a symmetric perturbation $Z$, differentiation of the trace of the matrix square root yields
\[
D\!\left[\tr R_s(X)\right][Z]
=
\frac12\tr\bigl(T_s(X)Z\bigr).
\]
\eqllead{Consequently,}{eq-gaussian-barycenter-first-variation}{
D\mathcal E(X)[Z]
=
\tr\bigl((\Id-\overline T(X))Z\bigr).
}
The same identity for singular $\cov_s$ follows by replacing $\cov_s$ with $\cov_s+\delta\Id$ and letting $\delta\downarrow0$. Since the minimizer lies in $\mathbb S_{++}^d$, first-order optimality is $\overline T(\cov)=\Id$. Finally, definitions~\eqref{eq-gaussian-barycenter-maps}--\eqref{eq-gaussian-barycenter-average-map} give
\(\Psi_\lambda(X) = X^{1/2}\overline T(X)X^{1/2},\)
so $\overline T(\cov)=\Id$ is equivalent to $\Psi_\lambda(\cov)=\cov$. Conversely, a positive-definite solution of this equation is stationary for the convex objective~\eqref{eq-gaussian-barycenter-energy}, hence is its unique minimizer. In dimension one, $\Bb(\sigma^2,\sigma_s^2)^2=(\sigma-\sigma_s)^2$, whose minimizer is $\sigma=\sum_s\lambda_s\sigma_s$.
\end{proof}
\(x\longmapsto \mean_s+T_s(X)(x-\mean).\)
\eql{\label{eq-gaussian-barycenter-transport-update}
\mathcal K_1(X)
\eqdef
\overline T(X)X\overline T(X)
=
X^{-1/2}\Psi_\lambda(X)^2X^{-1/2}.
}
\(B_\eta(X) \eqdef (1-\eta)\Id+\eta\overline T(X)\)
\begin{prop}[Convergence of the Gaussian transport fixed-point iteration]\label{prop-gaussian-barycenter-fixed-point-convergence}
Assume that $\lambda_s>0$, $\cov_s\succeq0$, and that at least one $\cov_s$ is positive definite. Fix $\eta\in(0,1]$ and $X^{(0)}\succ0$, and iterate
\(X^{(\ell+1)}=\mathcal K_\eta(X^{(\ell)}).\)
Then every $X^{(\ell)}$ is positive definite and $X^{(\ell)}\to\cov$, where $\cov$ is the covariance of the Gaussian barycenter in Proposition~\ref{prop-gaussian-barycenter}. More precisely, with
\(\mathcal R(X) \eqdef \tr\bigl((\Id-\overline T(X))X(\Id-\overline T(X))\bigr),\)
\eqllead{one has}{eq-gaussian-barycenter-descent}{
\mathcal E(X^{(\ell+1)})
\leq
\mathcal E(X^{(\ell)})
-
\eta(2-\eta)\mathcal R(X^{(\ell)}),
}
and therefore $\sum_{\ell\geq0}\mathcal R(X^{(\ell)})<+\infty$. Equivalently, \(\Bb\bigl(X^{(\ell)},X^{(\ell+1)}\bigr)^2 = \eta^2\mathcal R(X^{(\ell)}).\)
\end{prop}
\begin{proof}
Fix $X\succ0$, abbreviate $T_s=T_s(X)$, $\overline T=\overline T(X)$ and $B=B_\eta(X)$, and let $Z\sim\Gaussian(0,X)$. Since $T_sZ\sim\Gaussian(0,\cov_s)$, the pair $(BZ,T_sZ)$ is an admissible coupling between $\Gaussian(0,\mathcal K_\eta(X))$ and $\Gaussian(0,\cov_s)$. For every $z\in\RR^d$, the Euclidean variance decomposition around $\overline Tz=\sum_s\lambda_sT_sz$ gives
\begin{align*}
\sum_s\lambda_s\norm{Bz-T_sz}^2
&=
\sum_s\lambda_s\norm{\overline Tz-T_sz}^2
+(1-\eta)^2\norm{z-\overline Tz}^2,
\\
\sum_s\lambda_s\norm{z-T_sz}^2
&=
\sum_s\lambda_s\norm{\overline Tz-T_sz}^2
+\norm{z-\overline Tz}^2.
\end{align*}
Integrating and bounding each Wasserstein cost by the admissible coupling $(BZ,T_sZ)$ proves~\eqref{eq-gaussian-barycenter-descent}. Moreover, $B\succ0$: this is immediate when $\eta<1$, and for $\eta=1$ it follows because the weighted sum $\overline T$ contains at least one positive-definite term. Hence $z\mapsto Bz$ is the optimal map from $\Gaussian(0,X)$ to $\Gaussian(0,\mathcal K_\eta(X))$, which proves the displayed identity for the Bures step length.

Set
\(\delta(X)\eqdef\det(X)^{1/(2d)}, \qquad c_0\eqdef\sum_s\lambda_s\det(\cov_s)^{1/(2d)}>0.\)
The Minkowski determinant inequality, applied to the positive-semidefinite matrices $\Id$ and $T_s(X)$ with weights $1-\eta$ and $\eta\lambda_s$, gives
\begin{align*}
\delta\bigl(\mathcal K_\eta(X)\bigr)
&=
\delta(X)\det(B_\eta(X))^{1/d}
\\
&\geq
(1-\eta)\delta(X)
+
\eta\sum_s\lambda_s
\delta(X)\det(T_s(X))^{1/d}
\\
&=
(1-\eta)\delta(X)+\eta c_0,
\end{align*}
where the last identity uses $T_s(X)XT_s(X)=\cov_s$. It follows that
$\delta(X^{(\ell)})\geq\min\{\delta(X^{(0)}),c_0\}>0$ for every $\ell$.

The decreasing energy also bounds the traces. Indeed, the Schatten Cauchy--Schwarz inequality in the Bures formula gives \(\Bb(X,\cov_s)^2 \geq \pa{\sqrt{\tr X}-\sqrt{\tr\cov_s}}^2,\)
so~\eqref{eq-gaussian-barycenter-descent} bounds $\tr X^{(\ell)}$ uniformly. A uniform trace bound together with the determinant lower bound places all iterates in a compact subset of $\mathbb S_{++}^d$.

Summing~\eqref{eq-gaussian-barycenter-descent} gives $\mathcal R(X^{(\ell)})\to0$. Every subsequence therefore has a further subsequence converging to some $X_\infty\succ0$. Continuity of the matrix square root and inverse on the preceding compact set yields $\mathcal R(X_\infty)=0$, hence $\overline T(X_\infty)=\Id$. Proposition~\ref{prop-gaussian-barycenter} identifies $X_\infty$ with the unique covariance $\cov$. Since every cluster point is $\cov$, the full sequence converges to $\cov$. The undamped case $\eta=1$ is the Gaussian specialization of Theorem~4.2 of \'Alvarez-Esteban et al.; the argument above also proves the damped extension.
\end{proof}
\paragraph{Sinkhorn for barycenters.}
\eql{\label{eq-entropic-bary}
\umin{\a\in\simplex_n}
\sum_{s=1}^S\la_s\MKD_{\C_s}^\epsilon(\a,\b_s)
}
\eql{\label{eq-bary-entropy-couplings}
\umin{(\P_s)_s}
\enscond{
\epsilon\sum_s\la_s\KLD(\P_s|\K_s)
}{
\foralls s,\ \transp{\P_s}\ones_n=\b_s,
\quad
\P_1\ones_{n_1}=\ldots=\P_S\ones_{n_S}
},
}
where $\K_s\eqdef e^{-\C_s/\epsilon}$. The barycenter $\a$ is implicitly encoded in the common row marginal

\(\a=\P_1\ones_{n_1}=\ldots=\P_S\ones_{n_S}.\)
\eql{\label{eq-bary-opt}
\P_s=\diag(\uD_s)\K_s\diag(\vD_s),
}
\begin{align}
\foralls s\in\range{1,S},\qquad
\itt{\vD}_s
&\eqdef
\frac{\b_s}{\transp{\K_s}\it{\uD}_s},
\label{eq-sinkhorn-bary}
\\
\foralls s\in\range{1,S},\qquad
\itt{\uD}_s
&\eqdef
\frac{\itt{\a}}{\K_s\itt{\vD}_s},
\label{eq-sinkhorn-bary-2}
\\
\qwhereq
\itt{\a}
&\eqdef
\prod_s(\K_s\itt{\vD}_s)^{\la_s}.
\label{eq-sinkhorn-bary-3}
\end{align}
\(\gD_s^+ = \epsilon\log\pa{\frac{\b_s}{\transp{\K_s}e^{\fD_s/\epsilon}}},\)
\(\q_s \eqdef \K_s e^{\gD_s^+/\epsilon}, \qquad \a^+ = \prod_s\q_s^{\lambda_s}, \qquad \fD_s^+ = \epsilon\log\pa{\frac{\a^+}{\q_s}}.\)
Thus a complete generalized Sinkhorn cycle is exact two-block coordinate ascent on the dual, or equivalently alternating minimization of its negative. Indeed, the constraint follows from $\log\a^+=\sum_s\lambda_s\log\q_s$, while the first-order conditions require $e^{\fD_s^+/\epsilon}\q_s=\a^+$ for every $s$.

\begin{prop}[Dual of entropic barycenters]\label{prop-dual-entropic-barycenters}
Under the preceding positivity assumptions, the optimal scalings in~\eqref{eq-bary-opt} can be written as $(\uD_s,\vD_s)=(e^{\fD_s/\epsilon},e^{\gD_s/\epsilon})$, where $(\fD_s,\gD_s)_s$ solve the dual problem
\eql{\label{eq-dual-bary-entropy}
\umax{(\fD_s,\gD_s)_s}
\enscond{
\sum_s\la_s
\left(
\dotp{\gD_s}{\b_s}
-\epsilon\dotp{\K_s e^{\gD_s/\epsilon}}{e^{\fD_s/\epsilon}}
+\epsilon\sum_{i,j}(\K_s)_{i,j}
\right)
}{
\sum_s\la_s\fD_s=0
}.
}
\end{prop}
\begin{proof}
Introduce Lagrange multipliers in~\eqref{eq-bary-entropy-couplings}:
\begin{align*}
\umin{(\P_s)_s,\a}
\umax{(\fD_s,\gD_s)_s}
\sum_s\la_s\Big(
\epsilon\KLD(\P_s|\K_s)
+\dotp{\a-\P_s\ones_{n_s}}{\fD_s}
+\dotp{\b_s-\transp{\P_s}\ones_n}{\gD_s}
\Big).
\end{align*}
The explicit constraint $\a\in\simplex_n$ may be dropped here: nonnegativity of the couplings, together with $\P_s\ones_{n_s}=\a$ and $\transp{\P_s}\ones_n=\b_s$, already forces $\a\in\simplex_n$. We may therefore minimize the Lagrangian over $\a\in\RR^n$. Strong duality holds, so one can exchange the minimum and maximum. Finiteness of the minimum with respect to $\a$ gives the vector constraint $\sum_s\la_s\fD_s=0$, while minimization with respect to $\P_s$ gives the Legendre transform of $\KLD(\cdot|\K_s)$:
\[
\umax{(\fD_s,\gD_s)_s}
\sum_s\la_s
\left[
\dotp{\gD_s}{\b_s}
-\epsilon
\KLD^*\left(\frac{\fD_s\oplus\gD_s}{\epsilon}\middle|\K_s\right)
\right],
\qquad
\sum_s\la_s\fD_s=0.
\]
\eqllead{The separable conjugate is}{eq-legendre-kl-bary}{
\KLD^*(\VectMode{U}|\K)
=
\sum_{i,j}\K_{i,j}\big(e^{\VectMode{U}_{i,j}}-1\big),
}
because for $k>0$, \(\sup_{r\geq0} ur-\big(r\log(r/k)-r+k\big) = k(e^u-1),\)
and the case $k=0$ follows by lower semicontinuity. Substituting this conjugate gives~\eqref{eq-dual-bary-entropy}, including its additive constant. The coordinate maximization in $\gD_s$ gives~\eqref{eq-sinkhorn-bary}; the block maximization in all $(\fD_s)_s$ under $\sum_s\lambda_s\fD_s=0$ gives the weighted geometric mean~\eqref{eq-sinkhorn-bary-3} and then~\eqref{eq-sinkhorn-bary-2}.
\end{proof}
\paragraph{Wasserstein-over-Wasserstein and barycenters.}
\(\int_{\Pp(\Omega)}\MK_c(\beta_0,\alpha)\d\mathfrak A(\alpha)<+\infty,\)
\(\mathcal B_c(\mathfrak A) \eqdef \operatorname*{Argmin}_{\beta\in\Pp(\Omega)} \int_{\Pp(\Omega)}\MK_c(\beta,\alpha)\d\mathfrak A(\alpha).\)
\(\widetilde\alpha_{\mathfrak A} \eqdef \uargmin{\beta\in\Pp(\Omega)} \int_{\Pp(\Omega)}\MK_c(\beta,\alpha)\d\mathfrak A(\alpha),\)
\begin{prop}[Law of large numbers for barycenters over measures]\label{prop-wow-barycenter-lln}
Let $\Omega$ be a compact metric space, let $c:\Omega\times\Omega\to\RR$ be continuous, and let $\mathfrak A\in\Pp(\Pp(\Omega))$. Let $\alpha_1,\alpha_2,\ldots$ be independent random probability measures with common law $\mathfrak A$, and set
\(\widehat{\mathfrak A}_p \eqdef \frac1p\sum_{i=1}^p\delta_{\alpha_i} \in\Pp(\Pp(\Omega)).\)
Then, almost surely,
\begin{equation}\label{eq-wow-empirical-law-lln}
\widehat{\mathfrak A}_p \rightharpoonup \mathfrak A\ \text{ in }\Pp(\Pp(\Omega)),
\qquad
\bar\alpha_{\widehat{\mathfrak A}_p}\rightharpoonup\bar\alpha_{\mathfrak A}\ \text{ in }\Pp(\Omega).
\end{equation}
Moreover, if $\mathcal B_c(\mathfrak A)=\{\widetilde\alpha_{\mathfrak A}\}$ is a singleton and if $\widetilde\alpha_{\widehat{\mathfrak A}_p}\in\mathcal B_c(\widehat{\mathfrak A}_p)$ is any empirical barycenter, then, almost surely,
\eql{\label{eq-wow-population-barycenter-lln}
\widetilde\alpha_{\widehat{\mathfrak A}_p}
\rightharpoonup
\widetilde\alpha_{\mathfrak A}
\quad\text{in }\Pp(\Omega).
}
\end{prop}
\begin{proof}
Set $K=\Pp(\Omega)$. Since $\Omega$ is compact metric, $K$ is compact metric for weak convergence, and so is $\Pp(K)$. The space $\Cc(K)$ is separable for the uniform norm. Applying the scalar strong law to a countable dense family of test functions and then using uniform approximation gives, almost surely, for every $\Phi\in\Cc(K)$,
\(\int \Phi(\alpha)\d\widehat{\mathfrak A}_p(\alpha) = \frac1p\sum_{i=1}^p\Phi(\alpha_i) \longrightarrow \int \Phi(\alpha)\d\mathfrak A(\alpha)\)
This is exactly~\eqref{eq-wow-empirical-law-lln}.

For the collapsed mixtures, take $f\in\Cc(\Omega)$ and define $\Phi_f(\alpha)=\int_\Omega f\d\alpha$. This function is continuous on $\Pp(\Omega)$. Hence
\[
\int_\Omega f\d\bar\alpha_{\widehat{\mathfrak A}_p}
=
\int_{\Pp(\Omega)}\Phi_f(\alpha)\d\widehat{\mathfrak A}_p(\alpha)
\longrightarrow
\int_{\Pp(\Omega)}\Phi_f(\alpha)\d\mathfrak A(\alpha)
=
\int_\Omega f\d\bar\alpha_{\mathfrak A},
\]
which proves the second convergence in~\eqref{eq-wow-empirical-law-lln}.

Since $c$ is continuous on the compact product $\Omega\times\Omega$, the value $(\beta,\alpha)\mapsto\MK_c(\beta,\alpha)$ is continuous on $K^2$. Therefore the map $\beta\mapsto h_\beta$, where $h_\beta(\alpha)=\MK_c(\beta,\alpha)$, is continuous from the compact space $K$ to $\Cc(K)$. Its image
\(\mathcal H \eqdef \{h_\beta:\beta\in K\}\)
is compact in $\Cc(K)$. The convergence of $\widehat{\mathfrak A}_p$ to $\mathfrak A$ is uniform over $\mathcal H$: given $\eta>0$, cover $\mathcal H$ by finitely many $\eta$-balls in $\norm{\cdot}_\infty$, use weak convergence for the centers, and bound the error on the balls by the total variation of $\widehat{\mathfrak A}_p-\mathfrak A$. Hence the empirical objectives
\(F_p(\beta)\eqdef\int\MK_c(\beta,\alpha)\d\widehat{\mathfrak A}_p(\alpha)\)
converge uniformly on $K$ to
\(F(\beta)\eqdef\int\MK_c(\beta,\alpha)\d\mathfrak A(\alpha).\)
Let $\beta_p\in\mathcal B_c(\widehat{\mathfrak A}_p)$. Compactness gives a subsequence $\beta_{p_k}\rightharpoonup\beta$. Uniform convergence, continuity of $F$, and optimality of $\beta_{p_k}$ give, for any $\gamma\in K$,
\(F(\beta) = \lim_k F_{p_k}(\beta_{p_k}) \leq \lim_k F_{p_k}(\gamma) = F(\gamma),\)
so $\beta\in\mathcal B_c(\mathfrak A)$. If this set is the singleton $\{\widetilde\alpha_{\mathfrak A}\}$, every converging subsequence has the same limit, and therefore the whole sequence $\widetilde\alpha_{\widehat{\mathfrak A}_p}$ converges to $\widetilde\alpha_{\mathfrak A}$, proving~\eqref{eq-wow-population-barycenter-lln}.
\end{proof}
\paragraph{Toward central limit theorems on Wasserstein space.}
There are nevertheless important settings where such results can be proved. In one dimension, the quantile representation linearizes $\Wass_2$, so barycenter fluctuations can be studied through empirical averages of quantile functions.

\paragraph{Sliced and Radon barycenters.}
Slicing gives a scalable surrogate for high-dimensional barycenters by applying the one-dimensional quantile formula of Proposition~\ref{prop-quantile-barycenters} in every projection direction. For measures on $\RR^d$, define the sliced barycenter by replacing $\Wass_2$ with $\SW_2$, introduced in Definition~\ref{def-sliced-wasserstein} and interpreted through the Radon transform in Section~\ref{sec-sliced-wasserstein}:

\[
\umin{\al\in\Mm_+^1(\RR^d)}
\sum_{s=1}^S\la_s\SW_2(\al,\be_s)^2
=
\umin{\al\in\Mm_+^1(\RR^d)}
\int_{\Sphere^{d-1}}
\sum_{s=1}^S\la_s
\Wass_2\big((P_\theta)_\sharp\al,(P_\theta)_\sharp\be_s\big)^2
\d\sigma(\theta).
\]
\[
\umin{(\gamma_\theta)_\theta}
\int_{\Sphere^{d-1}}
\sum_{s=1}^S\la_s
\Wass_2(\gamma_\theta,\be_{s,\theta})^2
\d\sigma(\theta),
\qquad
\be_{s,\theta}\eqdef(P_\theta)_\sharp\be_s.
\]
\(\cumul{\gamma_\theta}^{-1}(r) = \sum_{s=1}^S\la_s\cumul{\be_{s,\theta}}^{-1}(r), \qquad r\in[0,1].\)
For two inputs $\be_0,\be_1$ and interpolation weight $t\in[0,1]$, it is useful to make this directionwise construction explicit. Define

\begin{equation}\label{eq-radon-barycenter-quantile-field}
Q_i(\theta,r)
\eqdef
\cumul{\be_{i,\theta}}^{-1}(r),
\qquad
Q_t(\theta,r)
\eqdef
(1-t)Q_0(\theta,r)+tQ_1(\theta,r),
\qquad
\gamma_{t,\theta}
\eqdef
\bigl(Q_t(\theta,\cdot)\bigr)_\sharp\mathrm{Leb}_{[0,1]}.
\end{equation}
Thus $Q_t$ is the quantile field of the independently interpolated projected laws $(\gamma_{t,\theta})_\theta$. When these laws have densities, we denote them by $h_t(\theta,\cdot)$.

Let $h(\theta,t)$ denote a density of $\gamma_\theta$. We use the one-dimensional Fourier transform in the variable $t$,

\begin{equation}\label{eq-radon-sinogram-fourier}
\widehat h(\theta,\omega)
\eqdef
\int_{\RR}e^{-\imath\omega t}h(\theta,t)\d t,
\qquad
h(\theta,t)
=
\frac{1}{2\pi}\int_{\RR}e^{\imath\omega t}\widehat h(\theta,\omega)\d\omega,
\end{equation}
\begin{prop}[Radon least-squares pseudoinverse]\label{prop-radon-pseudoinverse}
Let $d\geq2$, let $\sigma$ be the uniform probability measure on $\Sphere^{d-1}$, and let $R$ be the density Radon transform introduced in the \hyperref[rem-sliced-radon-viewpoint]{Radon point of view} paragraph. Write
\(\mathcal D(R) \eqdef \enscond{\rho\in L^2(\RR^d)}{R\rho\in L^2(\Sphere^{d-1}\times\RR,\d\sigma\,\d t)}.\)
Let $h\in L^2(\Sphere^{d-1}\times\RR,\d\sigma\,\d t)$ and assume that the Fourier expressions below define an element of $\mathcal D(R)$. Then the unique solution of
\begin{equation}\label{eq-radon-least-squares}
\umin{\rho\in\mathcal D(R)}
\int_{\Sphere^{d-1}}\int_{\RR}
\abs{R\rho(\theta,t)-h(\theta,t)}^2
\d t\,\d\sigma(\theta)
\end{equation}
is $\rho^\dagger=R^\dagger h$, where
\begin{equation}\label{eq-radon-pseudoinverse}
R^\dagger h(x)
=
\frac{\abs{\Sphere^{d-1}}}{2(2\pi)^d}
\int_{\Sphere^{d-1}}\int_{\RR}
e^{\imath\omega\dotp{\theta}{x}}
|\omega|^{d-1}
\widehat h(\theta,\omega)
\d\omega\,\d\sigma(\theta).
\end{equation}
Thus $R^\dagger h$ is the density of the pseudoinverse reconstruction, which is a priori a signed measure. If $h=R\rho$ is Radon-consistent and $\rho$ is sufficiently regular, then $R^\dagger R\rho=\rho$.
\end{prop}
\begin{proof}
For the $d$-dimensional Fourier transform, use the compatible convention
\[
\widehat\rho(\xi)
=
\int_{\RR^d}e^{-\imath\dotp{\xi}{x}}\rho(x)\d x,
\qquad
\rho(x)
=
\frac{1}{(2\pi)^d}\int_{\RR^d}e^{\imath\dotp{\xi}{x}}\widehat\rho(\xi)\d\xi.
\]
\eqllead{The Fourier-slice theorem gives}{eq-radon-fourier-slice}{
\widehat{R\rho}(\theta,\omega)=\widehat\rho(\omega\theta).
}
Applying the one-dimensional Plancherel identity in $t$ rewrites the objective in~\eqref{eq-radon-least-squares}, up to the positive factor $1/(2\pi)$, as
\(\int_{\Sphere^{d-1}}\int_{\RR} \abs{\widehat\rho(\omega\theta)-\widehat h(\theta,\omega)}^2 \d\omega\,\d\sigma(\theta).\)
Every $\xi\neq0$ has the two signed-polar representations
$(\theta,\omega)=(\xi/\norm{\xi},\norm{\xi})$ and
$(-\xi/\norm{\xi},-\norm{\xi})$. The pointwise least-squares minimizer is therefore
\begin{equation}\label{eq-radon-pseudoinverse-fourier}
\widehat{\rho^\dagger}(\xi)
=
\frac12\left(
\widehat h\left(\frac{\xi}{\norm{\xi}},\norm{\xi}\right)
+
\widehat h\left(-\frac{\xi}{\norm{\xi}},-\norm{\xi}\right)
\right).
\end{equation}
This also proves uniqueness under the stated existence assumptions. Applying the inverse $d$-dimensional Fourier transform to~\eqref{eq-radon-pseudoinverse-fourier} and using signed polar coordinates yields~\eqref{eq-radon-pseudoinverse}; the factor $\abs{\Sphere^{d-1}}/2$ accounts for the normalization of $\sigma$ and the two signed representations of each nonzero frequency. Finally, if $h=R\rho$, then~\eqref{eq-radon-fourier-slice} makes both terms in~\eqref{eq-radon-pseudoinverse-fourier} equal to $\widehat\rho(\xi)$, proving $R^\dagger R\rho=\rho$.
\end{proof}
\begin{equation}\label{eq-radon-windowed-ramp}
m_\Omega(\omega)
\eqdef
|\omega|^{d-1}\chi(\omega/\Omega),
\qquad
R_\Omega^\dagger h(x)
\eqdef
\frac{\abs{\Sphere^{d-1}}}{2(2\pi)^d}
\int_{\Sphere^{d-1}}\int_{\RR}
e^{\imath\omega\dotp{\theta}{x}}
m_\Omega(\omega)\widehat h(\theta,\omega)
\d\omega\,\d\sigma(\theta),
\end{equation}
where $\chi$ is an even low-pass window with $\chi(0)=1$. For the two-input path~\eqref{eq-radon-barycenter-quantile-field}, choose $\eta_t\geq0$ so that the positive part below has nonzero mass, and define the nonnegative, unit-mass reconstruction

\begin{equation}\label{eq-radon-display-reconstruction}
A_t(x)
\eqdef
\frac{\bigl((R_\Omega^\dagger h_t)(x)-\eta_t\bigr)_+}
{\displaystyle\int_{\RR^d}\bigl((R_\Omega^\dagger h_t)(z)-\eta_t\bigr)_+\d z}.
\end{equation}
For the endpoints, set $A_i=\rho_i$ when $\be_i=\rho_i\d x$, $i\in\{0,1\}$. The resulting regularized density is generally only a least-squares approximation to the independently averaged slices. With finitely many directions, the sampled Radon operator is also non-injective, so the pseudoinverse reconstruction and the constrained sliced barycenter may differ.

\subsection{Multimarginal OT}
\label{sec-multimarginal-ot}
\label{subsec:multi-marginal}
\paragraph{Definition and basic structure.}
\begin{defn}[Multimarginal optimal transport]\label{def-multimarginal-ot}
Let $\al_s\in\Pp(\Xx_s)$ for $s=1,\ldots,S$, and denote by
$\Couplings(\al_1,\ldots,\al_S)$ the probability measures whose
$s$-th marginal is $\al_s$. For a measurable cost $c$ bounded from below,
define
\begin{equation}\label{eq-multimarginal-ot}
\operatorname{MMOT}_c(\al_1,\ldots,\al_S)
\eqdef
\inf_{\pi\in\Couplings(\al_1,\ldots,\al_S)}
\int_{\Xx_1\times\cdots\times\Xx_S}
c(x_1,\ldots,x_S)\d\pi(x_1,\ldots,x_S).
\end{equation}
\end{defn}
\paragraph{Monge structure and splitting-set twist.}
\begin{defn}[Twist on splitting sets]\label{def-twist-splitting-sets}
Fix $x_1\in\Xx_1$. A set $M\subset\Xx_2\times\cdots\times\Xx_S$ is a $c$-splitting set at $x_1$ if there exist functions $u_s:\Xx_s\to\RR\cup\{-\infty\}$, for $s=2,\ldots,S$, such that
\(\sum_{s=2}^S u_s(x_s)\leq c(x_1,x_2,\ldots,x_S)\)
for all $(x_2,\ldots,x_S)$, with equality on $M$. Assume $c$ is differentiable in $x_1$. The cost is twisted on splitting sets if, for every $x_1$ and every $c$-splitting set $M$ at $x_1$, the map
\((x_2,\ldots,x_S) \longmapsto \nabla_{x_1}c(x_1,x_2,\ldots,x_S)\)
is injective on $M$.
\end{defn}
\begin{prop}[Multi-marginal Monge structure]\label{prop-multimarginal-monge-structure}
Assume that $\Xx_s\subset\RR^d$, that $c$ is continuous and differentiable with respect to $x_1$, and that $c$ is twisted on splitting sets. Suppose that Kantorovich dual maximizers $(\varphi_s)_{s=1}^S$ exist, that $\al_1$ is absolutely continuous, and that $\varphi_1$ is differentiable $\al_1$-a.e. Then every optimal plan $\pi^\star\in\Couplings(\al_1,\ldots,\al_S)$ is concentrated on the graph of maps
\(\pi^\star=(\Id,\T_2,\ldots,\T_S)_\sharp\al_1, \qquad (\T_s)_\sharp\al_1=\al_s.\)
In particular, under these hypotheses the optimizer is unique.
\end{prop}
\begin{proof}
Let $(\varphi_s)_{s=1}^S$ be optimal dual potentials. Complementary slackness gives a Borel contact set $\Gamma$ of full $\pi^\star$-measure on which $\sum_s\varphi_s(x_s)=c(x_1,\ldots,x_S)$. After disintegrating with respect to the first marginal, fix a point $x_1$ where $\varphi_1$ is differentiable and where the conditional plan is concentrated on the fiber
\(M(x_1)=\{(x_2,\ldots,x_S):(x_1,x_2,\ldots,x_S)\in\Gamma\}.\)
Indeed, set $u_2=\varphi_2+\varphi_1(x_1)$ and $u_s=\varphi_s$ for $s\geq3$. Dual feasibility gives $\sum_{s=2}^S u_s(x_s)\leq c(x_1,x_2,\ldots,x_S)$, and equality holds on $M(x_1)$. If $(x_2,\ldots,x_S)\in M(x_1)$, the function
\(z\longmapsto c(z,x_2,\ldots,x_S)-\sum_{s=2}^S\varphi_s(x_s)\)
touches $\varphi_1$ from above at $z=x_1$. Differentiating at this contact point gives \(\nabla\varphi_1(x_1)=\nabla_{x_1}c(x_1,x_2,\ldots,x_S).\)
All points in the fiber therefore have the same value of $\nabla_{x_1}c$. Twist on splitting sets makes the fiber a singleton for $\al_1$-a.e. $x_1$. Disintegrating $\pi^\star$ with respect to its first marginal gives Dirac conditional measures, hence measurable maps $(\T_2,\ldots,\T_S)$. If $\pi^1$ and $\pi^2$ are two optimal plans, their average is also optimal. The conditional measure of this average over $x_1$ is the average of the two Dirac conditionals, and it must again be a Dirac mass by the preceding argument. Hence the two Dirac masses coincide for $\al_1$-a.e. $x_1$, proving uniqueness.
\end{proof}
\paragraph{Coulomb cost and density-functional theory.}
\(c_{\mathrm{Coul}}(x_1,\ldots,x_N) \eqdef \sum_{1\leq i<j\leq N}\frac{1}{\norm{x_i-x_j}},\)
\(\nabla_{x_1}c_{\mathrm{Coul}}(x_1,\ldots,x_N) = -\sum_{j=2}^N\frac{x_1-x_j}{\norm{x_1-x_j}^3}.\)
For $N=2$, the map from $x_2$ to this vector is injective, so the ordinary two-marginal twist argument can be recovered under the required existence, duality and differentiability hypotheses. For $N\geq3$, however, the displayed total force does not by itself determine the entire tuple $(x_2,\ldots,x_N)$; twist on splitting sets, and hence a Monge representation, is not automatic.

If $\rho$ is an electron density with $\int_{\RR^3}\rho(x)\d x=N$ and $\al=\rho/N$ is the associated probability density, the strictly-correlated-electrons relaxation of density-functional theory is the equal-marginal problem

\[
V_{\mathrm{ee}}^{\mathrm{SCE}}[\rho]
\eqdef
\inf_{\pi\in\Couplings(\al,\ldots,\al)}
\int_{(\RR^3)^N}
c_{\mathrm{Coul}}(x_1,\ldots,x_N)
\d\pi(x_1,\ldots,x_N).
\]
\(\pi=(\Id,\T_2,\ldots,\T_N)_\sharp\al, \qquad (\T_i)_\sharp\al=\al,\)
\begin{prop}[Cyclic co-motion plans]\label{prop-cyclic-comotion-plans}
Let $\al\in\Pp(\RR^3)$ and let $\T:\RR^3\to\RR^3$ be measurable with $\T_\sharp\al=\al$ and $\T^N=\Id$ $\al$-a.e. Set $\T^0=\Id$ and
\(\pi_\T=(\T^0,\T^1,\ldots,\T^{N-1})_\sharp\al.\)
Then $\pi_\T\in\Couplings(\al,\ldots,\al)$. If $R_\sigma(x_1,\ldots,x_N)=(x_{\sigma(1)},\ldots,x_{\sigma(N)})$ for $\sigma\in\Perm(N)$, then
\(\bar\pi_\T\eqdef\frac1{N!}\sum_{\sigma\in\Perm(N)}(R_\sigma)_\sharp\pi_\T\)
is a symmetric admissible plan and
\[
\int c_{\mathrm{Coul}}\d\bar\pi_\T
=
\int c_{\mathrm{Coul}}\d\pi_\T
=
\int_{\RR^3}
\sum_{0\leq i<j\leq N-1}
\frac{1}{\norm{\T^i(x)-\T^j(x)}}\d\al(x).
\]
\end{prop}
\begin{proof}
Since $\T_\sharp\al=\al$, all iterates $\T^i$ preserve $\al$, so every marginal of $\pi_\T$ is $\al$. The plan $\bar\pi_\T$ is an average of coordinate permutations of $\pi_\T$, hence has the same marginals and is invariant under coordinate permutations. The Coulomb cost is symmetric in its arguments, so its integral is unchanged by each $R_\sigma$. The last identity follows by evaluating $c_{\mathrm{Coul}}$ on the graph $(x,\T(x),\ldots,\T^{N-1}(x))$.
\end{proof}
Proposition~\ref{prop-multimarginal-monge-structure} explains the general mechanism that can force graph solutions, while Proposition~\ref{prop-cyclic-comotion-plans} records the additional equal-marginal symmetry used by co-motion maps. Thus the DFT problem is both a central application and a warning that multi-marginal OT is richer than a naive deterministic matching problem.

\paragraph{Multi-marginal formulation of barycenters.}
\(c_{\mathrm{bar}}(x_1,\ldots,x_S) = \min_{x\in\RR^d} \sum_{s=1}^S\la_s\norm{x-x_s}^2.\)
\begin{prop}[Multi-marginal formula for quadratic barycenters]\label{prop-multimarginal-barycenter}
Let $\be_1,\ldots,\be_S\in\Pp_2(\RR^d)$ and $\la\in\simplex_S$. Define
\(B(x_1,\ldots,x_S)=\sum_{s=1}^S\la_s x_s, \qquad c_{\mathrm{bar}}(x_1,\ldots,x_S)=\min_x\sum_s\la_s\norm{x-x_s}^2.\)
If $\pi^\star$ solves the multi-marginal OT problem with marginals $(\be_s)_s$ and cost $c_{\mathrm{bar}}$, then $\al^\star=B_\sharp\pi^\star$ is a Wasserstein barycenter. Conversely, every barycenter is obtained this way from an optimal multi-marginal plan.
\end{prop}
\begin{proof}
For any candidate barycenter $\al$ and couplings $\pi_s\in\Couplings(\al,\be_s)$, glue the couplings along their common $\al$ marginal to obtain a joint law of $(X,Y_1,\ldots,Y_S)$. Conditioning on $(Y_s)_s$ and minimizing over $X$ gives
\(\sum_s\la_s\EE\norm{X-Y_s}^2 \geq \EE\min_x\sum_s\la_s\norm{x-Y_s}^2 = \EE c_{\mathrm{bar}}(Y_1,\ldots,Y_S).\)
Taking the infimum over the couplings gives that the barycenter value is at least the multi-marginal value. Conversely, from an optimal multi-marginal plan $\pi^\star$, set $X=B(Y_1,\ldots,Y_S)$. The couplings between $X$ and each $Y_s$ are feasible for the barycenter problem and attain exactly the multi-marginal cost, proving equality and the formula.
If $\al^\star$ is any barycenter, choose optimal couplings between $\al^\star$ and each $\be_s$ and glue them along the common $\al^\star$ marginal. Since the barycenter and multi-marginal values are equal, the conditional minimization inequality above must be an equality. Thus $X=B(Y_1,\ldots,Y_S)$ almost surely for the induced optimal multi-marginal plan, and $\al^\star=B_\sharp\pi^\star$.
\end{proof}
\begin{cor}[Sparse discrete barycenters]\label{cor-discrete-barycenters}
Let
\(\be_s=\sum_{i_s=1}^{n_s}b_{s,i_s}\de_{x_{s,i_s}}, \qquad s=1,\ldots,S,\)
be discrete input measures. There exists a quadratic Wasserstein barycenter $\al^\star$ supported on weighted averages $\sum_s\la_s x_{s,i_s}$ of one support point from each input and satisfying
\(\#\supp(\al^\star) \leq \sum_{s=1}^S n_s-S+1.\)
\end{cor}
\begin{proof}
Write a discrete multi-marginal plan as a nonnegative tensor
$\P=(\P_{i_1,\ldots,i_S})$ and let $\mathcal A_{\mathrm{marg}}$ collect its $S$ marginals: \(\bigl(\mathcal A_{\mathrm{marg}}\P\bigr)_{s,i_s} = \sum_{(i_r)_{r\neq s}}\P_{i_1,\ldots,i_S}.\)
After vectorizing $\P$, Proposition~\ref{prop-lp-rank-sparsity} applies to the multi-marginal linear program. Its constraint operator has rank
\(\operatorname{rank}(\mathcal A_{\mathrm{marg}}) = \sum_{s=1}^S n_s-S+1.\)
Indeed, a family of vectors $u_s\in\RR^{n_s}$ belongs to the annihilator of its image precisely when
\(\sum_{s=1}^S u_{s,i_s}=0 \qquad\text{for every }(i_1,\ldots,i_S).\)
Varying one index at a time shows that every $u_s$ is constant, say $u_s=c_s\ones_{n_s}$, and the remaining condition is $\sum_s c_s=0$. The annihilator therefore has dimension $S-1$, which proves the rank formula.

The rank-controlled sparsity result therefore supplies an optimal multi-marginal tensor $\P^\star$ with at most $\sum_s n_s-S+1$ positive entries. Proposition~\ref{prop-multimarginal-barycenter} gives the barycenter $\al^\star=B_\sharp\P^\star$. Each positive tensor entry produces at most one atom under $B$, and collisions can only reduce the support size. This proves the claim.
\end{proof}
\paragraph{Entropic regularization of multi-marginal OT.}
\(\inf_{\pi\in\Couplings(\al_1,\ldots,\al_S)} \int c\d\pi+\epsilon\KL(\pi|\al_1\otimes\cdots\otimes\al_S).\)
\[
\d\pi^\star(x_1,\ldots,x_S)
=
\exp\!\left(\frac{\sum_s f_s(x_s)-c(x_1,\ldots,x_S)}{\epsilon}\right)
\prod_s\d\al_s(x_s),
\]
\paragraph{Treewidth and graphical structure.}
The important exception is when the cost factors over a sparse interaction graph. Let $G=(V,E)$ be a finite undirected graph with $V=\{1,\ldots,S\}$ and suppose, for simplicity, that

\begin{equation}\label{eq-multimarginal-graph-cost}
c(x_1,\ldots,x_S)
=
\sum_{(r,s)\in E}c_{r,s}(x_r,x_s).
\end{equation}
\begin{defn}[Tree decomposition and treewidth]\label{def-treewidth}
A \emph{tree decomposition} of a finite graph $G=(V,E)$ is a tree $\mathcal T=(Q,F)$ together with bags $B_q\subseteq V$, $q\in Q$, such that
\begin{enumerate}
\item $\bigcup_{q\in Q}B_q=V$;
\item for every edge $(r,s)\in E$, some bag contains both $r$ and $s$;
\item for every vertex $s\in V$, the nodes $\enscond{q\in Q}{s\in B_q}$ form a connected subtree of $\mathcal T$.
\end{enumerate}
The width of this decomposition is $\max_{q\in Q}|B_q|-1$, and the \emph{treewidth} $\operatorname{tw}(G)$ is the minimum width over all tree decompositions.
\end{defn}
\paragraph{Junction-tree contractions inside Sinkhorn.}
\[
\K^{r,s}_{i_r,i_s}
=
\exp\!\left(-\frac{\C^{r,s}_{i_r,i_s}}{\epsilon}\right),
\qquad
h_s(i_s)
=
(\a_s)_{i_s}(u_s)_{i_s}.
\]
\begin{equation}\label{eq-multimarginal-graph-factorization}
\P_{i_1,\ldots,i_S}
=
\prod_{s\in V}h_s(i_s)
\prod_{(r,s)\in E}\K^{r,s}_{i_r,i_s}.
\end{equation}
When $G$ is a tree, let $N_G(r)$ denote the neighbors of $r$. The directed sum-product messages satisfy

\begin{equation}\label{eq-multimarginal-tree-message}
m_{r\to s}(i_s)
=
\sum_{i_r=1}^{n_r}
\K^{r,s}_{i_r,i_s}h_r(i_r)
\prod_{q\in N_G(r)\setminus\{s\}}
m_{q\to r}(i_r).
\end{equation}
\begin{equation}\label{eq-multimarginal-tree-marginal}
(\widehat{\a}_s)_{i_s}
=
h_s(i_s)
\prod_{r\in N_G(s)}m_{r\to s}(i_s).
\end{equation}
For the block $s$ currently selected by the cyclic scaling algorithm, the exact coordinate update is then

\[
u_s
\leftarrow
u_s\odot\frac{\a_s}{\widehat{\a}_s},
\]
with componentwise products and quotients. This changes only the unary factor $h_s$.

For a general tree decomposition, choose one host bag for each unary factor, assign each edge factor in~\eqref{eq-multimarginal-graph-factorization} to a bag containing both endpoints, and denote the product assigned to bag $B_q$ by $\psi_q(i_{B_q})$. For adjacent bags $q,q'\in Q$, set $\mathcal S_{q,q'}=B_q\cap B_{q'}$. Junction-tree messages take the form

\begin{equation}\label{eq-multimarginal-junction-message}
M_{q\to q'}(i_{\mathcal S_{q,q'}})
=
\sum_{i_{B_q\setminus\mathcal S_{q,q'}}}
\psi_q(i_{B_q})
\prod_{\ell\in N_{\mathcal T}(q)\setminus\{q'\}}
M_{\ell\to q}(i_{\mathcal S_{\ell,q}}).
\end{equation}
\begin{equation}\label{eq-multimarginal-junction-belief}
\mathfrak b_q(i_{B_q})
=
\psi_q(i_{B_q})
\prod_{\ell\in N_{\mathcal T}(q)}
M_{\ell\to q}(i_{\mathcal S_{\ell,q}}).
\end{equation}
For any $s\in B_q$, summing $\mathfrak b_q$ over $B_q\setminus\{s\}$ gives $\widehat{\a}_s$; the running-intersection property ensures that every bag containing $s$ gives the same result.

Redundant adjacent bags may be contracted, so one can assume that the decomposition is reduced. If its width is $w$, every bag contains at most $w+1$ variables and every separator contains at most $w$.

\[
O\!\left(|Q|\,n^{w+1}\right)
\quad\text{time},
\qquad
O\!\left(|Q|\,n^w\right)
\quad\text{message memory}.
\]
\[
O\!\left(
\sum_{q\in Q}
\bigl(1+\deg_{\mathcal T}(q)\bigr)
\prod_{s\in B_q}n_s
\right),
\]
\[
O\!\left(
\sum_{(q,q')\in F}
\prod_{s\in\mathcal S_{q,q'}}n_s
\right)
\]
\subsection{Low-Rank Optimal Transport}
\label{sec-low-rank-ot}
\begin{defn}[Low-rank factored couplings]\label{def-low-rank-couplings}
Let $\a\in\simplex_n$, $\b\in\simplex_m$ and let $r\geq1$. A rank-$r$ factored coupling is a triple $(\Q,\R,g)$ such that
\(\Q\in\RR_+^{n\times r},\qquad \R\in\RR_+^{m\times r}, \qquad g\in\simplex_r,\)
with \(\Q\ones_r=\a, \qquad \R\ones_r=\b, \qquad \transp \Q\ones_n=\transp \R\ones_m=g.\)
It induces the coupling
\begin{equation}\label{eq-low-rank-coupling-factor}
\P(\Q,\R,g)
=
\Q\diag(g)^{-1}\transp \R,
\qquad
\P_{i,j}=\sum_{k=1}^r \frac{\Q_{i,k}\R_{j,k}}{g_k},
\end{equation}
where columns with $g_k=0$ are discarded before applying the formula.
\end{defn}
\begin{prop}[Factored couplings and nonnegative rank]\label{prop-low-rank-factorization}
For every feasible triple $(\Q,\R,g)$ in Definition~\ref{def-low-rank-couplings}, the matrix $\P(\Q,\R,g)$ belongs to $\CouplingsD(\a,\b)$ and has nonnegative rank at most $r$. Conversely, every coupling $\P\in\CouplingsD(\a,\b)$ with nonnegative rank at most $r$ admits a representation of the form~\eqref{eq-low-rank-coupling-factor}, after possibly deleting zero latent components.
\end{prop}
\begin{proof}
The marginal constraints follow by direct summation: \(\P\ones_m=\Q\diag(g)^{-1}\transp \R\ones_m =\Q\diag(g)^{-1}g=\Q\ones_r=\a,\)
and similarly $\transp \P\ones_n=\b$. Since $\P$ is a sum of $r$ nonnegative rank-one matrices $\Q_{:,k}\R_{:,k}^\top/g_k$, its nonnegative rank is at most $r$.

Conversely, suppose $\P=UV^\top$ with $U\in\RR_+^{n\times r}$ and $V\in\RR_+^{m\times r}$. Set $q_k=\sum_i U_{i,k}$ and $s_k=\sum_j V_{j,k}$. Components with $q_ks_k=0$ do not contribute and can be removed. Define $g_k=q_ks_k$, $\Q_{i,k}=U_{i,k}s_k$ and $\R_{j,k}=V_{j,k}q_k$. Since $\P$ has total mass one, $g\in\simplex_r$. Moreover $\Q\ones_r=\P\ones_m=\a$, $\R\ones_r=\transp \P\ones_n=\b$, and both column marginals are equal to $g$. Finally,
\(\sum_k\frac{\Q_{i,k}\R_{j,k}}{g_k} = \sum_k\frac{U_{i,k}s_kV_{j,k}q_k}{q_ks_k} =\P_{i,j}.\)
\end{proof}
For a cost matrix $\C\in\RR^{n\times m}$, the low-rank constrained OT value is therefore

\(\min_{(\Q,\R,g)} \dotp{\C}{\Q\diag(g)^{-1}\transp \R} \quad\text{subject to Definition~\ref{def-low-rank-couplings}.}\)
\begin{equation}\label{eq-low-rank-entropic-ot}
\min_{\substack{\Q\ones_r=\a,\ \Q^\top\ones_n=g\\
\R\ones_r=\b,\ \R^\top\ones_m=g}}
\dotp{\C}{\Q\diag(g)^{-1}\transp \R}
+
\epsilon\KLD(\Q|\a\otimes g)
+
\epsilon\KLD(\R|\b\otimes g),
\end{equation}
where $\Q,\R$ are nonnegative. For fixed $g$, this differs only by constants from adding $-\epsilon(\HD(\Q)+\HD(\R))$ to the factorized transport cost.

\subsection{Capacity-Constrained Optimal Transport}
\label{sec-capacity-constrained-ot}
\begin{defn}[Capacity-constrained optimal transport]\label{def-capacity-constrained-ot}
Let $\alpha\in\Mm_+^1(\Xx)$, $\beta\in\Mm_+^1(\Yy)$, let $c:\Xx\times\Yy\to[0,+\infty]$ be measurable, and let $u:\Xx\times\Yy\to[0,+\infty]$ be a measurable capacity. Define
\begin{equation}\label{eq-capacity-constrained-ot}
\MK_c^u(\alpha,\beta)
\eqdef
\inf_{\pi\in\Couplings(\alpha,\beta)}
\enscond{
\int_{\Xx\times\Yy} c(x,y)\d\pi(x,y)
}{
\pi\ll\alpha\otimes\beta,\quad
\frac{\d\pi}{\d(\alpha\otimes\beta)}(x,y)\leq u(x,y)
\quad (\alpha\otimes\beta)\text{-a.e.}
}.
\end{equation}
By convention, $u\equiv+\infty$ removes both the density bound and the absolute-continuity requirement, so that $\MK_c^{+\infty}(\alpha,\beta)=\MK_c(\alpha,\beta)$.
\end{defn}
The product coupling $\alpha\otimes\beta$ is feasible whenever $u\geq1$. Thus the constraint is not meant to promote the product plan; rather, it prevents the optimizer from using any pair more than the prescribed density ratio.

For discrete measures $\alpha=\sum_i a_i\de_{x_i}$ and $\beta=\sum_j b_j\de_{y_j}$, a capacity is an upper matrix $\overline{\P}\in\RR_+^{n\times m}$. Given a cost matrix $\C$, the density-ratio discretization of Definition~\ref{def-capacity-constrained-ot}, with $\overline{\P}_{i,j}=u_{i,j}a_i b_j$, is the capped linear program

\begin{equation}\label{eq-discrete-capacity-constrained-ot}
\min_{\P\in\CouplingsD(\a,\b)}
\dotp{\C}{\P}
\quad\text{subject to}\quad
0\leq \P_{i,j}\leq \overline{\P}_{i,j}\quad\forall(i,j).
\end{equation}
\begin{prop}[Feasibility of a capped transport polytope]\label{prop-capacity-feasibility}
The constraints in~\eqref{eq-discrete-capacity-constrained-ot} are feasible if and only if
\begin{equation}\label{eq-capacity-cut-condition}
a(I)+b(J)-1\leq \overline{\P}(I,J)
\qquad
\text{for every }I\subset\{1,\ldots,n\},\ J\subset\{1,\ldots,m\}.
\end{equation}
\end{prop}
\begin{proof}
If $\P$ is feasible, then
\(\P(I,J) = a(I)-\P(I,J^c) \geq a(I)-b(J^c) = a(I)+b(J)-1.\)
Since $\P(I,J)\leq \overline{\P}(I,J)$, condition~\eqref{eq-capacity-cut-condition} is necessary.

For sufficiency, build a flow network with capacity $a_i$ from the source to row vertex $i$, capacity $\overline{\P}_{i,j}$ from row $i$ to column $j$, and capacity $b_j$ from column $j$ to the sink. A cut containing row set $I$ and column set $J^c$ has capacity
\(1-a(I)+\overline{\P}(I,J)+1-b(J)\).
Condition~\eqref{eq-capacity-cut-condition} says that every such cut has capacity at least one. The max-flow/min-cut theorem therefore gives a unit flow, and its row-to-column edge values form the required matrix $\P$.
\end{proof}
\begin{equation}\label{eq-entropic-capacity-constrained-ot}
\min_{\P\in\CouplingsD(\a,\b),\,0\leq \P\leq \overline{\P}}
\dotp{\C}{\P}
+
\epsilon\KLD(\P|\a\otimes\b).
\end{equation}
On the active support, all iterates and correction factors are positive, so every division in the KL--Dykstra iteration is well defined. KL-Dykstra convergence on finite-dimensional affine and box constraints then shows that, when the capped polytope is nonempty, the iterates converge to the unique minimizer of~\eqref{eq-entropic-capacity-constrained-ot}.

\subsection{Metric Learning and Inverse OT}
\label{sec-metric-learning-inverse-ot}
\subsubsection{Differentiating OT losses}
\paragraph{First variations of OT values.}
\begin{prop}[First variations of OT values]\label{prop-ot-first-variations-unregularized}\label{prop-ot-first-variations-entropic}
Let $\alpha$ and $\beta$ be probability measures with compact supports $\Xx=\supp(\alpha)$ and $\Yy=\supp(\beta)$, let $c\in\Cc(\Xx\times\Yy)$, and let $\epsilon\geq0$. Use the convention $\MK_c^0=\MK_c$, where the unregularized and positive-temperature values are defined in Definitions~\ref{def-continuous-kantorovich-problem} and~\ref{def-continuous-entropic-ot}. Denote by $\mathcal O_c^\epsilon(\alpha,\beta)$ and $\mathcal D_c^\epsilon(\alpha,\beta)$ the sets of primal and dual optimizers of the corresponding problem.
If $\chi$ is a signed measure with $\chi(\Xx)=0$ and $\alpha_t=\alpha+t\chi$ is a probability measure for $0\leq t\leq t_0$, then
\[
\left.\frac{\d}{\d t}\right|_{t=0^+}
\MK_c^\epsilon(\alpha_t,\beta)
=
\sup_{(f,g)\in\mathcal D_c^\epsilon(\alpha,\beta)}
\int_{\Xx} f(x)\d\chi(x).
\]
If $h\in\Cc(\Xx\times\Yy)$ and $c_t=c+th$, then
\[
\left.\frac{\d}{\d t}\right|_{t=0^+}
\MK_{c_t}^\epsilon(\alpha,\beta)
=
\inf_{\pi\in\mathcal O_c^\epsilon(\alpha,\beta)}
\int_{\Xx\times\Yy} h(x,y)\d\pi(x,y).
\]
For fixed $(\alpha,\beta)$, a finite signed Radon measure $\eta$ on $\Xx\times\Yy$ represents the first variation of the functional $c\mapsto\MK_c^\epsilon(\alpha,\beta)$ if, for every $h\in\Cc(\Xx\times\Yy)$,
\(\MK_{c+t h}^\epsilon(\alpha,\beta) = \MK_c^\epsilon(\alpha,\beta) +t\int_{\Xx\times\Yy}h\d\eta+o(t) \qquad(t\to0).\)
In particular, if the normalized optimal potential $f_\epsilon^\star$ and the optimal plan $\pi_\epsilon^\star$ are unique (so in particular when $\epsilon>0$), then, with the marginal first variation understood as in Definition~\ref{def-first-variation},
\[
\frac{\delta \MK_c^\epsilon}{\delta\alpha}(\alpha,\beta)=f_\epsilon^\star,
\qquad
\frac{\delta \MK_c^\epsilon}{\delta c}(\alpha,\beta)=\pi_\epsilon^\star.
\]
\end{prop}
\begin{proof}
For $\epsilon=0$, the Kantorovich duality theorem is stated in Proposition~\ref{prop-kantorovich-duality-general}. Its dual value is a supremum of affine functions of $\alpha$, so Danskin's theorem gives the first formula as the supremum over active dual potentials. The condition $\chi(\Xx)=0$ makes this expression independent of the additive gauge $(f,g)\mapsto(f+\lambda,g-\lambda)$.

For $\epsilon>0$, use the continuous entropic duality and primal--dual density law of Proposition~\ref{prop-continuous-entropic-duality}. At an optimal pair, the marginal constraint makes the density in~\eqref{eq-continuous-entropic-density-law} integrate to one with respect to $\beta$ for every $x$ on the source support. Therefore the derivative of the dual objective~\eqref{eq-dual-sinkhorn-objective} with respect to $\alpha$ at fixed optimal potentials is $\int f_\epsilon^\star\d\chi$: the derivative of its exponential term vanishes after the row normalization. Danskin's theorem then gives the same first formula, now with a singleton set of normalized active potentials.

For every $\epsilon\geq0$, the primal objective depends on $c$ only through the affine term $\int c\d\pi$. The minimum form of Danskin's theorem therefore gives the second formula as the infimum of $\int h\d\pi$ over active primal optimizers. When the normalized potential or plan is unique, the corresponding directional derivative is linear in the perturbation, which gives the two first variations.
\end{proof}
\(f^\star\in\partial_{\a}\MKD_{\C}(\a,\b), \qquad \P^\star\in\partial_{\C}^{\mathrm{sup}}\MKD_{\C}(\a,\b).\)
For $\epsilon>0$, these are ordinary derivatives on the relative interiors of finite-dimensional simplices. For a finite-dimensional parametrization $c_\theta$, the entropic formula gives the backpropagation rule

\(\partial_{\theta}\MK_{c_\theta}^\epsilon(\alpha,\beta) = \int \partial_\theta c_\theta(x,y)\d\pi_{\theta,\epsilon}^\star(x,y),\)
where $\pi_{\theta,\epsilon}^\star$ is the entropic optimizer between $(\alpha,\beta)$ for the cost $c_\theta$. For every perturbation $H\in\RR^{d\times d}$,

\[
D_A\MK_{c_A}^\epsilon(\alpha,\beta)[H]
=
\int \dotp{Hx}{y}\d\pi_{A,\epsilon}^\star(x,y)
=
\left\langle H,\int yx^\top\d\pi_{A,\epsilon}^\star(x,y)\right\rangle_{\mathrm F}.
\]
\[
\nabla_A\MK_{c_A}^\epsilon(\alpha,\beta)
=
\int yx^\top\d\pi_{A,\epsilon}^\star(x,y)
=
\operatorname{Cov}_{\pi_{A,\epsilon}^\star}(Y,X)+m_\beta m_\alpha^\top,
\]
where $(X,Y)\sim\pi_{A,\epsilon}^\star$, $m_\alpha=\int x\d\alpha(x)$ and $m_\beta=\int y\d\beta(y)$. Thus the gradient is the raw cross-moment of the optimal coupling, and it is its cross-covariance when both marginals are centered.

\subsubsection{Inverse Optimal Transport}
\begin{defn}[Regularized inverse-OT loss]\label{def-inverse-ot-loss}
Let $\epsilon\geq0$ and $\widehat\pi\in\Couplings(\alpha,\beta)$, with $\KL(\widehat\pi|\alpha\otimes\beta)<+\infty$ when $\epsilon>0$. The regularized inverse-OT loss is
\eql{\label{eq-inverse-ot-loss}
\mathcal F_\epsilon(c\mid\widehat\pi)
\eqdef
\int c\d\widehat\pi
+\epsilon\KL(\widehat\pi|\alpha\otimes\beta)
-\MK_c^\epsilon(\alpha,\beta),
}
where the entropy term is set to zero when $\epsilon=0$.
\end{defn}
As shown next, this loss is convex in $c$. Minimizing it over a convex class of costs is therefore a convex problem that avoids non-convex bilevel optimization through a forward OT solver.
\begin{prop}[Convexity and calibration of the inverse-OT loss]\label{prop-inverse-ot-convex}
For fixed $(\alpha,\beta,\widehat\pi)$, the map $c\mapsto\mathcal F_\epsilon(c\mid\widehat\pi)$ is convex and nonnegative. Moreover,
\(\mathcal F_\epsilon(c\mid\widehat\pi)=0 \quad\Longleftrightarrow\quad \widehat\pi\text{ solves~\eqref{eq-entropic-generic} for the cost }c.\)
When $\epsilon>0$, the regularized optimizer is unique whenever the value is finite, so the zero set identifies the forward coupling for a fixed cost.
\end{prop}
\begin{proof}
The first two terms in~\eqref{eq-inverse-ot-loss} are respectively linear and constant in $c$, whereas $c\mapsto\MK_c^\epsilon(\alpha,\beta)$ is concave because it is the infimum of affine functions of $c$, as in Proposition~\ref{prop-kantorovich-value-curvature}. Testing~\eqref{eq-entropic-generic} with $\widehat\pi$ proves nonnegativity, and equality holds precisely when this test plan attains the infimum. Uniqueness for $\epsilon>0$ follows from the strict convexity of relative entropy on finite-entropy couplings.
\end{proof}
\paragraph{Bilinear cost learning.}
For concreteness, consider the bilinear model on $\RR^d$,

\(c_A(x,y)=\dotp{Ax}{y}, \qquad A\in\RR^{d\times d}.\)
\eql{\label{eq-inverse-ot-bilinear-estimator}
A^\star\in\uargmin{A\in\mathcal A}
\mathcal F_\epsilon(c_A\mid\widehat\pi).
}
\paragraph{Finite-sample and population geometry.}
Let $A_0=-\Id$, let $\pi_0$ be the associated quadratic $\Wass_2$ plan, and set $\widehat\pi_n=n^{-1}\sum_{i=1}^n\de_{(X_i,Y_i)}$ for i.i.d. $(X_i,Y_i)\sim\pi_0$. Along any affine matrix path $A_t$, the empirical loss $t\mapsto\mathcal F_0(c_{A_t}\mid\widehat\pi_n)$ is convex and polyhedral. Under the smoothness, invertibility and Hessian-spanning assumptions of Peyr\'e, Poon and Tron (2026), its population limit instead has local quadratic curvature transverse to the positive-scaling ray of $A_0$.

\paragraph{Statistical estimation.}
\(\norm{A_n^\star-A_0}_{\mathrm F}=O_\PP(n^{-1/2})\)
\subsection{Weak Optimal Transport}
\label{sec-weak-ot}
\paragraph{Barycentric projection of a coupling.}
\begin{defn}[Barycentric projection of a coupling]\label{def-barycentric-projection}
Let $\alpha,\beta\in\Pp_1(\RR^d)$ and let $\pi\in\Couplings(\alpha,\beta)$. Disintegrate $\pi$ with respect to its first marginal as
$\pi(\d x,\d y)=\pi_x(\d y)\alpha(\d x)$. Since $\beta$ has finite first moment, the conditional mean is finite for $\alpha$-a.e. $x$, and the barycentric projection of $\pi$ is the map
\begin{equation}\label{eq-barycentric-projection}
\bar T_\pi(x)
\eqdef
\int_{\RR^d}y\d\pi_x(y),
\qquad
\bar\beta_\pi
\eqdef
(\bar T_\pi)_\sharp\alpha.
\end{equation}
\end{defn}
For discrete measures $\alpha=\sum_{i=1}^n a_i\de_{x_i}$ and $\beta=\sum_{j=1}^m b_j\de_{y_j}$, write $\pi=\sum_{i,j}\P_{i,j}\de_{(x_i,y_j)}$ with $\P\in\CouplingsD(\a,\b)$. Then
\(
\pi_{x_i}=\sum_{j=1}^m \frac{\P_{i,j}}{a_i}\de_{y_j}
\qandq
\bar T_\pi(x_i)=\frac{1}{a_i}\sum_{j=1}^m\P_{i,j}y_j.
\)
\begin{prop}[Barycentric projection of a quadratic optimal plan]\label{prop-barycentric-projection-optimal}
Let $\pi\in\Couplings(\alpha,\beta)$ be optimal for the quadratic cost $\norm{x-y}^2$ between $\alpha,\beta\in\Pp_2(\RR^d)$, and define $\bar T_\pi$ and $\bar\beta_\pi$ by~\eqref{eq-barycentric-projection}. Then $(\Id,\bar T_\pi)_\sharp\alpha$ is an optimal coupling between $\alpha$ and $\bar\beta_\pi$. Equivalently, $\bar T_\pi$ is a quadratic optimal transport map from $\alpha$ to the projected target $\bar\beta_\pi$.
\end{prop}
\begin{proof}
By Theorem~\ref{thm:opt_ccm}, $\pi$ is concentrated on a $c$-cyclically monotone set $\Gamma$ for $c(x,y)=\norm{x-y}^2$. For the quadratic cost, and since it is enough to check cyclic permutations, this means that every finite cycle $(x_i,y_i)_{i=1}^m\subset\Gamma$ satisfies
\(\sum_{i=1}^m \dotp{x_i}{y_i} \geq \sum_{i=1}^m \dotp{x_i}{y_{i+1}}, \qquad y_{m+1}=y_1.\)
After changing the disintegration on an $\alpha$-negligible set, $\pi_x$ is supported on the section
\(\Gamma_x=\enscond{y}{(x,y)\in\Gamma}\)
for $\alpha$-a.e. $x$.
Choose $x_1,\ldots,x_m$ in this full-measure set and independently sample $Y_i\sim\pi_{x_i}$. Applying the cyclic inequality to $(x_i,Y_i)$ and taking expectations gives
\(\sum_{i=1}^m \dotp{x_i}{\bar T_\pi(x_i)} \geq \sum_{i=1}^m \dotp{x_i}{\bar T_\pi(x_{i+1})}.\)
Thus $(\Id,\bar T_\pi)_\sharp\alpha$ is concentrated on a cyclically monotone graph. By the cyclic-monotonicity characterization of quadratic optimality, this plan is optimal between its two marginals, namely $\alpha$ and $\bar\beta_\pi$.
\end{proof}
\paragraph{Weak transport costs.}
\begin{defn}[Weak optimal transport]\label{def-weak-optimal-transport}
Let $\Xx$ and $\Yy$ be Polish spaces, let $\alpha\in\Pp(\Xx)$ and $\beta\in\Pp(\Yy)$, and let $C:\Xx\times\Pp(\Yy)\to\RR\cup\{+\infty\}$ be measurable. For each $\pi\in\Couplings(\alpha,\beta)$, let
$\pi(\d x,\d y)=\alpha(\d x)\pi_x(\d y)$ be a disintegration. Then
\begin{equation}\label{eq-weak-ot}
\WOT_C(\alpha,\beta)
\eqdef
\inf_{\pi\in\Couplings(\alpha,\beta)}
\int C(x,\pi_x)\d\alpha(x).
\end{equation}
\end{defn}
If $C(x,\cdot)$ is convex for $\alpha$-a.e. $x$, then the objective in~\eqref{eq-weak-ot} is convex in $\pi$: the conditional law of $(1-t)\pi^0+t\pi^1$ is $(1-t)\pi_x^0+t\pi_x^1$. Since $\Couplings(\alpha,\beta)$ is convex, weak OT is then a convex optimization problem; without this hypothesis it need not be.
\begin{prop}[Weak Kantorovich duality]\label{prop-weak-ot-duality}
Let $\Xx$ and $\Yy$ be compact metric spaces, and let $C:\Xx\times\Pp(\Yy)\to\RR\cup\{+\infty\}$ be proper, jointly lower semicontinuous, bounded from below and convex in its second argument. Fix $\alpha\in\Pp(\Xx)$ and $\beta\in\Pp(\Yy)$, and assume that $\WOT_C(\alpha,\beta)<+\infty$. For $g\in\Cc(\Yy)$, define the weak $C$-transform
\[
g^C(x)
\eqdef
\inf_{\nu\in\Pp(\Yy)}
\left\{
C(x,\nu)-\int g(y)\d\nu(y)
\right\}.
\]
Then
\[
\WOT_C(\alpha,\beta)
=
\sup_{g\in\Cc(\Yy)}
\left\{
\int g^C(x)\d\alpha(x)+\int g(y)\d\beta(y)
\right\}.
\]
When $C(x,\nu)=\int c(x,y)\d\nu(y)$, this reduces to the usual Kantorovich dual with $g^C(x)=\inf_y(c(x,y)-g(y))$.
\end{prop}
\begin{proof}
For any coupling $\pi$ and any $g\in\Cc(\Yy)$, the definition of $g^C$ gives
\(C(x,\pi_x)\geq g^C(x)+\int g(y)\d\pi_x(y)\).
After integration with respect to $\alpha$, the second term becomes $\int g\d\beta$ because the second marginal of $\pi$ is $\beta$. This proves weak duality.

For the converse, lift a kernel $x\mapsto\pi_x$ to the measure $\alpha(\d x)\delta_{\pi_x}(\d\nu)$ on $\Xx\times\Pp(\Yy)$. Relaxing this graph constraint gives probability measures $P$ whose first marginal is $\alpha$ and whose intensity satisfies
\(\int_{\Xx\times\Pp(\Yy)}\nu\,\d P(x,\nu)=\beta\).
This relaxation does not change the value: disintegrating $P(\d x,\d\nu)=P_x(\d\nu)\alpha(\d x)$ and replacing $P_x$ by its barycenter $\bar\nu_x=\int\nu\,\d P_x(\nu)$ preserves the intensity constraint and cannot increase the objective, by convexity of $C(x,\cdot)$. The relaxed feasible set is compact, and the integral functional is lower semicontinuous. Fenchel--Rockafellar duality for the affine intensity constraint therefore has no gap; its continuous multiplier is $g\in\Cc(\Yy)$. Minimization in $\nu$ gives $g^C(x)$, while the constraint contributes $\int g\d\beta$. This is the weak-cost analogue of Proposition~\ref{prop-duality-discr}; see for the Polish-space formulation.
\end{proof}
The dual objective is concave in $g$, because $g^C(x)$ is a pointwise infimum of affine functions of $g$; it is generally non-smooth.

The most common weak transport model on $\RR^d$ retains only the conditional barycenter.
\begin{prop}[Barycentric weak transport is weaker than $\Wass_2$]\label{prop-barycentric-weak-ot}
Let $\alpha,\beta\in\Pp_2(\RR^d)$ and define
\(C_{\mathrm{bar}}(x,\nu) = \norm{x-\int y\d\nu(y)}^2.\)
Equivalently,
\begin{equation}\label{eq-barycentric-weak-primal}
\WOT_{C_{\mathrm{bar}}}(\alpha,\beta)
=\inf_{\pi\in\Couplings(\alpha,\beta)}
\int\norm{x-\bar T_\pi(x)}^2\d\alpha(x).
\end{equation}
Then \(\WOT_{C_{\mathrm{bar}}}(\alpha,\beta) \leq \Wass_2^2(\alpha,\beta).\)
\end{prop}
\begin{proof}
Let $\pi$ be any coupling and disintegrate it as $\pi_x\alpha$. By Jensen's inequality, \(\norm{x-\bar T_\pi(x)}^2 \leq \int\norm{x-y}^2\d\pi_x(y).\)
Integrating in $x$ gives $\int C_{\mathrm{bar}}(x,\pi_x)\d\alpha(x)\leq\int\norm{x-y}^2\d\pi(x,y)$. Taking the infimum over $\pi$ proves the claim.
\end{proof}
For $\alpha=\sum_i a_i\de_{x_i}$ and $\beta=\sum_j b_j\de_{y_j}$, this is the convex quadratic program
\begin{equation}\label{eq-barycentric-weak-discrete}
\min_{\P\in\CouplingsD(\a,\b)}
\sum_i a_i\norm{x_i-\frac{1}{a_i}\sum_j\P_{i,j}y_j}^2.
\end{equation}
Unlike the generally non-convex quadratic GW problem of Section~\ref{sec-gromov-wasserstein}, every row barycenter here is affine in $\P$. Adding $\epsilon\KL(\pi|\alpha\otimes\beta)$ gives a convex entropic problem. If $F_{\mathrm{bar}}$ denotes the objective in~\eqref{eq-barycentric-weak-primal}, its linearized cost, modulo a source-only gauge, is
\begin{equation}\label{eq-barycentric-weak-linearized-cost}
c_\pi^{\mathrm{bar}}(x,y)\eqdef 2\dotp{x-\bar T_\pi(x)}{x-y},
\qquad
\left.\frac{\d}{\d t}\right|_{t=0}F_{\mathrm{bar}}((1-t)\pi+t\widehat\pi)
=\int c_\pi^{\mathrm{bar}}\d(\widehat\pi-\pi).
\end{equation}
Hence a regularized minimizer $\pi^\star$ solves the entropic OT problem with frozen cost $c_{\pi^\star}^{\mathrm{bar}}$, and convexity makes this condition globally sufficient. This suggests Sinkhorn subproblems as in the iterative GW solver of Section~\ref{sec-gromov-wasserstein}, but the tangent is a lower model, so a damped or proximal update is needed for a generic descent guarantee.
The barycentric cost is the canonical example to keep in mind: admissibility still constrains the full conditional laws to have second marginal $\beta$, but the objective only charges the displacement from $x$ to $\bar T_\pi(x)$ and ignores the conditional variance around this barycenter. Vanishing cost means $\bar T_\pi=\Id$, which is exactly the martingale constraint studied next.

\subsubsection{Martingale Optimal Transport}
\label{sec-martingale-ot}
\begin{defn}[Martingale couplings and martingale OT]\label{def-martingale-coupling}
Let $\alpha,\beta\in\Pp_1(\RR^d)$. A coupling $\pi\in\Couplings(\alpha,\beta)$ is a martingale coupling if, for the disintegration $\pi(\d x,\d y)=\pi_x(\d y)\alpha(\d x)$,
\begin{equation}\label{eq-martingale-coupling}
\int_{\RR^d}y\d\pi_x(y)=x
\qquad\text{for $\alpha$-a.e. }x.
\end{equation}
Equivalently, $\bar T_\pi=\Id$ in~\eqref{eq-barycentric-projection}. The set of such couplings is denoted
\(\Couplings_{\mathrm{mart}}(\alpha,\beta) \eqdef \enscond{\pi\in\Couplings(\alpha,\beta)}{\bar T_\pi=\Id\quad\alpha\text{-a.e.}}.\)
\end{defn}
\begin{rem}[Zero barycentric cost and martingale feasibility]\label{rem-weak-ot-martingale-feasibility}
For $\alpha,\beta\in\Pp_2(\RR^d)$,
\begin{equation}\label{eq-weak-zero-iff-martingale}
\WOT_{C_{\mathrm{bar}}}(\alpha,\beta)=0
\quad\Longleftrightarrow\quad
\Couplings_{\mathrm{mart}}(\alpha,\beta)\neq\emptyset.
\end{equation}
A martingale coupling is a zero-cost competitor. Conversely, existence for the barycentric weak problem and nonnegativity of~\eqref{eq-barycentric-weak-primal} show that a zero value yields an optimizer with $\bar T_\pi=\Id$. By Theorem~\ref{thm-strassen-martingale}, this is also equivalent to $\alpha\preceq_{\mathrm{cx}}\beta$.
\end{rem}
For a cost $c$, the martingale OT problem minimizes $\int c\d\pi$ over $\Couplings_{\mathrm{mart}}(\alpha,\beta)$, with value $+\infty$ when the admissible set is empty.

The terminology comes from probability: if $(X,Y)\sim\pi$, then~\eqref{eq-martingale-coupling} is exactly $\EE[Y|X]=X$. Hence martingale OT is a Kantorovich problem with the usual two marginal constraints plus a barycentric constraint on the conditional laws.

\paragraph{Stochastic orders.}
\begin{defn}[Classical stochastic order]\label{def-classical-stochastic-order}
For $\alpha,\beta\in\Pp(\RR)$,
\[
\alpha\preceq_{\mathrm{st}}\beta
\quad\Longleftrightarrow\quad
\int\varphi\,\d\alpha\leq\int\varphi\,\d\beta
\quad\text{for every bounded increasing }\varphi.
\]
\end{defn}
\begin{prop}[Strassen's theorem for stochastic order]\label{prop-strassen-stochastic-order}
Let $\alpha,\beta\in\Pp(\RR)$. Then $\alpha\preceq_{\mathrm{st}}\beta$ if and only if there exists $\pi\in\Couplings(\alpha,\beta)$ concentrated on $\enscond{(x,y)}{x\leq y}$.
\end{prop}
\begin{proof}
A coupling supported on $x\leq y$ immediately gives the integral inequality for every increasing $\varphi$. Conversely, testing approximate indicators of intervals $(t,+\infty)$ gives $F_\alpha(t)\geq F_\beta(t)$ for all continuity points of the distribution functions. If $U$ is uniform on $(0,1)$ and $F_\alpha^{-1},F_\beta^{-1}$ are the generalized quantiles, then $X=F_\alpha^{-1}(U)$ and $Y=F_\beta^{-1}(U)$ have laws $\alpha$ and $\beta$, and the inequality between distribution functions gives $X\leq Y$ almost surely.
\end{proof}
\paragraph{Convex order and martingale feasibility.}
\begin{defn}[Convex order]\label{def-convex-order}
For $\alpha,\beta\in\Pp_1(\RR^d)$, write $\alpha\preceq_{\mathrm{cx}}\beta$ when
\(\int \varphi\,\d\alpha\leq\int \varphi\,\d\beta\)
for every continuous convex $\varphi:\RR^d\to\RR$ with at most linear growth.
\end{defn}
\begin{thm}[Strassen's martingale theorem]\label{thm-strassen-martingale}
Let $\alpha,\beta\in\Pp_1(\RR^d)$. Then
\(\alpha\preceq_{\mathrm{cx}}\beta \quad\Longleftrightarrow\quad \Couplings_{\mathrm{mart}}(\alpha,\beta)\neq\emptyset.\)
Equivalently, $\beta$ is above $\alpha$ in convex order if and only if there exists a coupling $\pi(\d x,\d y)=\pi_x(\d y)\alpha(\d x)$ such that $\int y\d\pi_x(y)=x$ for $\alpha$-almost every $x$.
\end{thm}
\begin{proof}
If $\pi$ is a martingale coupling, then Jensen's inequality gives, for every convex $\varphi$,
\[
\varphi(x)
=
\varphi\!\left(\int y\d\pi_x(y)\right)
\leq
\int \varphi(y)\d\pi_x(y),
\]
and integration in $x$ gives $\int\varphi\d\alpha\leq\int\varphi\d\beta$.

Conversely, assume $\alpha\preceq_{\mathrm{cx}}\beta$. First suppose that both measures are supported in a compact convex set $K$. Separate $\beta$ from the closed convex set of probability measures on $K$ that are reachable from $\alpha$ by martingale kernels supported in $K$. If $\beta$ were not in this set, a continuous separating function $\psi:K\to\RR$ would satisfy
\[
\int\psi\d\beta
>
\sup
\left\{
\int\psi\d\eta
\;:\;
\eta\in\Pp(K),\
\Couplings_{\mathrm{mart}}(\alpha,\eta)\neq\emptyset
\right\}.
\]
For a fixed $x\in K$, optimizing over probability measures on $K$ with barycenter $x$ gives the concave envelope $\operatorname{conc}_K\psi(x)$; hence the right-hand side is $\int\operatorname{conc}_K\psi\d\alpha$. Since convex order is equivalently the reverse inequality for concave functions,
\(\int\operatorname{conc}_K\psi\d\alpha \geq \int\operatorname{conc}_K\psi\d\beta \geq \int\psi\d\beta,\)
a contradiction. For general measures in $\Pp_1(\RR^d)$, the same separation argument is carried out in the $\Wass_1$ topology, whose continuous test functions have at most linear growth. Compactness is replaced by tightness and uniform integrability of first moments; these properties also make the set of attainable second marginals closed and preserve the martingale constraint under limits. This is the standard extension in Strassen's theorem.
\end{proof}
\begin{prop}[Brenier--Strassen projection formula]\label{prop-brenier-strassen-projection}
Let $\alpha,\beta\in\Pp_2(\RR^d)$ and let $C_{\mathrm{bar}}$ be the barycentric quadratic cost of Proposition~\ref{prop-barycentric-weak-ot}. Then
\begin{equation}\label{eq-brenier-strassen-projection}
\WOT_{C_{\mathrm{bar}}}(\alpha,\beta)
=
\inf_{\eta\in\Pp_2(\RR^d):\,\eta\preceq_{\mathrm{cx}}\beta}
\Wass_2^2(\alpha,\eta).
\end{equation}
\end{prop}
\begin{proof}
Let $(X,Y)$ have law $\pi\in\Couplings(\alpha,\beta)$, set $Z=\EE[Y\mid X]$, and denote its law by $\eta$. Conditional Jensen gives
\(\EE[\varphi(Z)] \leq \EE[\varphi(Y)]\)
for every convex $\varphi$, so $\eta\preceq_{\mathrm{cx}}\beta$. Since $(X,Z)$ couples $\alpha$ and $\eta$, \(\Wass_2^2(\alpha,\eta) \leq \EE\norm{X-Z}^2 = \int C_{\mathrm{bar}}(x,\pi_x)\d\alpha(x).\)
Taking the infimum over $\pi$ proves that the left-hand side of~\eqref{eq-brenier-strassen-projection} is no smaller than the right-hand side.

Conversely, fix $\eta\preceq_{\mathrm{cx}}\beta$. By Theorem~\ref{thm-strassen-martingale}, there is a martingale coupling from $\eta$ to $\beta$. Glue it to an optimal coupling $(X,Z)$ between $\alpha$ and $\eta$, conditionally independently given $Z$. The resulting variables satisfy $\EE[Y\mid Z]=Z$, hence
\(\EE[Y\mid X]=\EE[Z\mid X]\).
Conditional Jensen therefore gives \(\EE\norm{X-\EE[Y\mid X]}^2 = \EE\norm{X-\EE[Z\mid X]}^2 \leq \EE\norm{X-Z}^2 = \Wass_2^2(\alpha,\eta).\)
Taking the infimum over $\eta$ proves the reverse inequality.
\end{proof}
\(\Gaussian(m,\Sigma_0)\preceq_{\mathrm{cx}}\Gaussian(m,\Sigma_1) \quad\Longleftrightarrow\quad \Sigma_1-\Sigma_0\succeq0.\)
